% !TEX root = ../main.tex
\chapter{Dealers and Liquid Security Markets}\label{ch:dealers}
\begin{paracol}{2}

\begin{sources}
\begin{tabularx}{\linewidth}{@{}l l L@{}}
\toprule
\textbf{Segment} & \textbf{Length} & \textbf{Content}\\
\midrule
\seg{10}{1}{L10.1 FT: Asymmetric Credit Growth in Europe} & 6:33 & IMF: European banks deleveraging; deposits flee from periphery to core.\\
\seg{10}{2}{L10.2 Market Liquidity, Dealers, and Inventories} & 7:10 & Market versus funding liquidity; the supermarket and its inventories.\\
\seg{10}{3}{L10.3 Two-sided Dealer Basics} & 6:31 & Two inventories; dealers as smoothers in price and time.\\
\seg{10}{4}{L10.4 Economics of the Dealer Function: the Treynor Model} & 11:43 & Position limits, outside spread, inside spread.\\
\seg{10}{5}{L10.5 Leveraged Dealer Basics} & 7:22 & Dealing with borrowed money; the Volcker rule and matched book.\\
\seg{10}{6}{L10.6 Real World Dealers} & 7:56 & Gross versus net exposure; inventories ``out in the world''.\\
\seg{10}{7}{L10.7 Arbitrage and the Assumption of Perfect Liquidity} & 9:41 & Dealers move liquidity across markets; perfect liquidity as a limit case; Fischer Black's ``Noise''.\\
\bottomrule
\end{tabularx}

\smallskip
\textbf{Files.} Transcripts \file{materials/transcripts/L10-P*.txt}; Mehrling's notes \mn{10}{1}.
\textbf{Reading.} Jack Treynor, ``The Economics of the Dealer Function'', \emph{Financial Analysts Journal} (1987) (not among the course files); Fischer Black, ``Noise'' (1986) (\file{materials/readings/L10-Noise.pdf}); IMF GFSR Box~2.3 and the FT article (\file{materials/readings/L10-*.pdf}).
Recommended: Larry Harris, \emph{Trading and Exchanges: Market Microstructure for Practitioners} (2003).
\end{sources}

\section*{Motivation}
Lecture~9 showed banks as dealers in money in Bagehot's world. This lecture builds the general model of a dealer, Treynor's, which the course uses from now on:
for security dealers here, for banks as dealers in money in lecture~11, for foreign exchange and swaps later. The question it answers: where does
\emph{market liquidity} come from? Answer: from dealers, who supply it for a profit, and in doing so move prices away from fundamental values. Perfect liquidity,
the assumption behind most of finance and general equilibrium theory, is shown to be a logically impossible limit.

% ------------------------------------------------------------------
\section{FT: asymmetric credit growth in Europe}\label{sec:dl-ft}

\begin{casestudy}[FT, 9 October 2012: ``IMF warns eurozone on capital flight'']
The IMF's \emph{Global Financial Stability Report} warns that unless eurozone officials step up their response, European banks would dump \$2.8 trillion of assets, more
than 7\% of their balance sheets, by the end of next year; banks in the periphery about 10\% \ts{10}{1}{0:07}. Mehrling notes that the FT mislabelled one of the
IMF's charts, ``an F for accuracy'': pay attention to details \ts{10}{1}{0:48}.
The article (by Claire Jones and Michael Steen) adds three points:
\begin{itemize}
  \item The IMF calls the result ``extreme fragmentation'' of euro-area funding markets. Its estimates are based on the 58 largest EU banks.
  \item With additional national measures, deleveraging would be \$2.3 trillion, 6\% of balance sheets.
  \item Draghi: ``You can't have a union when you have certain countries that are permanent creditors and a set of countries that are permanent debtors.''
\end{itemize}
The chart is from Box~2.3 of the IMF \emph{Global Financial Stability Report}, October 2012. Its Figure~2.3.3 shows euro-area bank credit to the non-bank private sector (Dec 2009 = 100) rising in the core and falling in the periphery.
\end{casestudy}
Depositors move money from banks in the periphery (Spain, Italy, Portugal, Ireland) to banks in the core, fearing bankruptcy or a break-up of the euro, in which
they would end up with a liability of, say, the Irish rather than the German central bank \ts{10}{1}{1:28}. This involves movements of reserves
among central banks (TARGET2 balances, left for later) \ts{10}{1}{2:29}. For the periphery banks, losing deposits means they must find other funding, impossible, or
shrink assets: they shed securities and loans and try to raise capital, and they make no new loans \ts{10}{1}{2:50}.
\begin{taccts}
\tacct[2.3cm]{Periphery bank}{$-$Securities\\$-$Loans}{$-$Deposits\\$+$Capital}\qquad
\tacct[2.3cm]{Core bank}{$+$Reserves / securities}{$+$Deposits}
\end{taccts}
For European banks as a whole the leverage ratio, total liabilities to capital, fell from about 25 to 23 (banks have 3--4\% capital) \ts{10}{1}{3:32}. The average
hides the asymmetry: bank credit to the private sector shrinks in the periphery and keeps growing in the core \ts{10}{1}{4:37}. The IMF adds that the ECB's
pledge to buy government debt lowered sovereign yields, but it is too early to tell whether it relieves deleveraging pressure \ts{10}{1}{6:01}. The question of the lecture:
how can a mere \emph{pledge} to buy lower bond yields? \ts{10}{1}{6:01}

% ------------------------------------------------------------------
\section{Market liquidity and inventories}\label{sec:dl-liquidity}

\begin{definition}[New concepts: funding liquidity, market liquidity]\label{def:dl-liq}
\emph{Funding liquidity} is the ability to raise money, to replace funds, the problem of the periphery banks \ts{10}{2}{0:28}. \emph{Market liquidity} is the ability to buy or
sell an asset quickly, in volume, without moving the price much; it is a property of the market for a particular asset \ts{10}{2}{0:49}. In Hicks's terms,
a liquid market has a \emph{continuous} price: it may fluctuate, but without jumps or ``air pockets''; an illiquid market trades in fits and starts, at discontinuous prices
\ts{10}{2}{1:51}.
\end{definition}

Supply and demand say nothing about the market structure behind such continuity: the price simply moves from one equilibrium to another, an abstraction from the
institutions that make it continuous, the dealers \ts{10}{2}{3:01}.

\paragraph{The supermarket.} Intuition from a dealer we know, the West Side Market: in the early morning it is empty and the shelves full; in the evening ``all of New York''
is in there and the shelves are depleted; the prices are the same morning and evening \ts{10}{2}{3:42}. Supply and demand are not lined up in time, yet the price
does not move. The secret is \emph{inventories}: the shelves absorb the flow of demand and are restocked at night, from inventories in New Jersey, the wholesale market, which in
turn are restocked from farms: ``dare I say a hierarchy of inventories'' \ts{10}{2}{4:43}. The supermarket is a \emph{one-sided} dealer: it sells
retail but does not buy from you (come in with a box of mangoes and they will ask whether you stole it); it buys wholesale. An investment bank underwriting an initial public
offering is a one-sided security dealer: it buys the whole issue and distributes it, holding inventory meanwhile \ts{10}{2}{6:07}. Even Japanese ``just-in-time'' production
does not eliminate inventories, the notes observe: it pushes them back before production, into the parts and components held by suppliers \mn{10}{2}.

% ------------------------------------------------------------------
\section{Two-sided dealers}\label{sec:dl-twosided}

A two-sided dealer trades a security for money and money for the security, either way; so it needs two inventories, of cash and of securities. A simple (hypothetical)
dealer with fixed capital allocates it between the two and lets them fluctuate: when the public sells securities to it, its cash falls and its securities rise, and vice versa
\ts{10}{3}{0:28}.
\begin{taccts}
\tacct[2.5cm]{Unleveraged dealer}{Cash\\Securities}{Capital}
\end{taccts}
With large enough inventories the price would not move at all \ts{10}{3}{2:12}.

\paragraph{Life without a dealer.} A buyer must find a seller: either raise the price offered (``I know it's worth 10, but I'll pay 12'') or wait. Time and price are the two
costs \ts{10}{3}{2:33}. The housing market is the example: a broker shows you houses, and you wait; knocking on the door of an owner who is not ready to
sell, you overpay; needing to sell quickly, you cut the price \ts{10}{3}{3:56}. Dealers are \emph{smoothers}: they can flatten the price and reduce the
waiting time to nothing, because they are always ready to take the other side \ts{10}{3}{3:36}. They do it to make money, so we need the economics of the dealer
function \ts{10}{3}{4:58}.

\begin{intuition}[A missing piece of microeconomics]
Price theory, and general equilibrium, assume perfect liquidity: a market-clearing price at every moment. This assumption must be false: dealers supply market liquidity, and not
for free; if liquidity were a free good, infinite in quantity, no dealer would supply it. ``The abstraction we have in price theory is an abstraction that's logically impossible for
liquid markets'' \ts{10}{3}{5:19}.
\end{intuition}

% ------------------------------------------------------------------
\section{The economics of the dealer function: Treynor's model}\label{sec:dl-treynor}

Treynor's diagram has the dealer's \emph{inventory} on the horizontal axis and the \emph{price} on the vertical axis (\cref{fig:dl-treynor}). It has three elements.

\begin{nfigure}
\begin{tikzpicture}[font=\scriptsize,x=1cm,y=1cm]
  \draw[-{Stealth[length=4pt]}] (-4.6,0) -- (4.8,0) node[right]{inventory};
  \draw[-{Stealth[length=4pt]}] (0,-0.2) -- (0,5.2) node[above]{price};
  \draw[dashed] (-4,0) -- (-4,4.9); \draw[dashed] (4,0) -- (4,4.9);
  \node[below,align=center] at (-4,0) {dealer's\\maximum short};
  \node[below,align=center] at (4,0) {dealer's\\maximum long};
  \node[below] at (0,0) {0};
  % outside spread
  \node[defcol,rotate=-22] at (-2.0,4.05) {dealer's ask};
  \node[intcol,rotate=-22] at (-1.4,2.65) {dealer's bid};
  \draw[very thick,thmcol] (-4.3,4.4) -- (4.3,4.4) node[right,align=left]{value-based\\ask};
  \draw[very thick,thmcol] (-4.3,0.8) -- (4.3,0.8) node[right,align=left]{value-based\\bid};
  % inside spread: ask and bid lines sloping down, reaching the outside spread at the limits
  \draw[very thick,defcol] (-4,4.4) -- (4,1.15);
  \draw[very thick,intcol] (-4,4.05) -- (4,0.8);
  \draw[dotted] (0,2.6) -- (-0.15,2.6) node[left]{$p^\ast$};
  \draw[{Stealth[length=3pt]}-{Stealth[length=3pt]}] (1.5,1.74) -- (1.5,2.13);
  \node[right] at (1.55,1.95) {inside spread};
  \draw[{Stealth[length=3pt]}-{Stealth[length=3pt]}] (-4.55,0.8) -- (-4.55,4.4);
  \node[left,align=right] at (-4.6,2.6) {outside\\spread};
  \node[align=center,font=\scriptsize\itshape] at (2.6,3.6) {market liquidity\\(dealer)};
  \node[align=center,font=\scriptsize\itshape] at (0,-1.35) {funding liquidity: how far from 0\\the dealer can finance its position};
\end{tikzpicture}
\caption{Treynor's model of a dealer \ts{10}{4}{0:13}. The dealer quotes an ask and a bid (inside spread) that fall as its inventory grows; the value-based traders' prices (outside
spread) bound the market; the position limits come from the dealer's funding. $p^\ast$, the midpoint at zero inventory, can be read as fundamental value
\ts{10}{7}{4:07}.}\label{fig:dl-treynor}
\end{nfigure}

\begin{enumerate}
  \item \textbf{Position limits.} The dealer's maximum short and maximum long positions. A long position gains if the price rises, a short position gains if it falls; at zero
        inventory the dealer is not exposed \ts{10}{4}{0:36}. The limits come from capital, but more generally from the ability to \emph{borrow} to fund the inventory: a
        banker who will lend up to some amount and no more, or the dealer's own risk appetite. They are soft in reality, hard in the model, and symmetric only for convenience
        \ts{10}{4}{1:21}. This is \emph{funding liquidity}.
  \item \textbf{The outside spread.} Behind the scenes are the fundamental forces of supply and demand and deep pockets: the \emph{value-based traders}, ``Warren Buffett in the
        wings'', who buy when the price is low enough (the value-based bid) and sell or short when it is high enough (the value-based ask) \ts{10}{4}{3:07}. They are
        the dealer's ``liquidity provider of last resort'': a dealer near its long limit can unload its inventory to them. Knowing they are there, the dealer is willing to let
        inventory build up, and so to supply liquidity \ts{10}{4}{4:08}. In a crisis the value-based bid goes to zero; dealers stop making markets
        and prices are in free fall \ts{10}{4}{5:11}.
  \item \textbf{The inside spread.} The dealer's own quotes, bid and ask, are what you see in the market when your broker quotes you Apple stock; the outside spread you see only if
        inventories push the price to it \ts{10}{4}{5:51}. Two things determine it \ts{10}{4}{6:34}:
        \begin{enumerate}[(a)]
          \item its \emph{width}: the volatility of the price; holding inventory of a volatile asset is risky, and a wider spread protects the dealer \ts{10}{4}{6:55};
          \item its \emph{slope}: adverse selection; the dealer trades with anyone, and the counterparty may know more. A growing inventory means many people are selling to it: the
                order flow is information, ``maybe they know something I don't know'', so the dealer buys more only at a better price, and lowers both quotes
                \ts{10}{4}{7:38}.
        \end{enumerate}
\end{enumerate}

\begin{definition}[New concepts: inside and outside spread, adverse selection]
The \emph{inside} (dealer) spread is the difference between the dealer's ask and bid; the \emph{outside} spread is the difference between the prices at which value-based investors
will sell and buy. \emph{Adverse selection} is the risk of trading with a better-informed counterparty: the trades that happen are disproportionately those that are bad for you.
\end{definition}

\begin{proposition}[Dealers turn funding liquidity into market liquidity]
Security dealers demand funding liquidity, borrowing to finance their inventories, and supply market liquidity, quoting prices at which anyone can trade \ts{10}{4}{10:02}.
\end{proposition}

The notes describe the liquidity the dealer supplies as an \emph{option to trade}: at the quoted prices, either way, if you want to, but it is up to you. Anyone who places a limit order
supplies liquidity in this sense, and amateur day traders often do, ``until they run out of money'' (the concept is developed in Larry Harris, \emph{Trading and Exchanges}, 2003)
\mn{10}{4}. In the notes' Treynor model the value-based traders' outside spread may be ``maybe 20\% below and 20\% above'' their estimate of fundamental value; when the dealer hits his
position limits, the value-based trader becomes market maker of last resort \mn{10}{4}.

\begin{keyformulas}
\textbf{A linear Treynor dealer} (a simple parametrisation of \cref{fig:dl-treynor}): with inventory $I\in[-I_{\max},I_{\max}]$,
\[
  \text{ask}(I) = p^\ast + \frac{s}{2} - \lambda I ,\qquad \text{bid}(I) = p^\ast - \frac{s}{2} - \lambda I ,
\]
with inside spread $s$ (rising with volatility) and slope $\lambda$ (rising with adverse selection), chosen so that the quotes reach the value-based ask and bid at the position
limits: $\lambda I_{\max} + s/2 = $ half the outside spread.
\end{keyformulas}

\begin{exercise}[The pledge]
Use the Treynor diagram to explain how a central bank's credible pledge to buy a government bond at a given price can lower its yield even before any purchase: which element of the diagram does the
pledge change?
\end{exercise}

% ------------------------------------------------------------------
\section{Leveraged dealers and the Volcker rule}\label{sec:dl-leverage}

Dealing is a very competitive business, with pressure to quote narrow spreads; one way to beat the competition is to trade on less capital, with borrowed money, so that the profit on
capital is multiplied by leverage \ts{10}{5}{0:28}. Leverage is the ratio of debt to equity; deleveraging, as in the FT story, reduces it \ts{10}{5}{1:10}.
Capital-funded dealers exist, but they are the value-based traders: ask Warren Buffett what his balance sheet looks like, ``billions of dollars of cash waiting to buy a railroad''; the
inside spread is made by leveraged dealers \ts{10}{5}{1:52}. In the extreme, fully borrowed: the ``inventory of cash'' is the banker who lends \ts{10}{5}{2:33}.
\begin{taccts}
\textit{long:}\quad\tacct[2.0cm]{Leveraged dealer}{Securities}{Loan from bank}\qquad
\textit{short:}\quad\tacct[2.0cm]{Leveraged dealer}{Cash}{Securities (owed)}
\end{taccts}
Making markets now means expanding and contracting the balance sheet itself, with only a thin layer of capital to absorb price changes \ts{10}{5}{3:14}. Mehrling knows
of no regulatory leverage limit for dealers; they are usually part of larger organisations with more capital \ts{10}{5}{4:16}.

\paragraph{The Volcker rule again.} It says that banks cannot have subsidiaries that take positions: they may deal, but with a \emph{matched book}, net inventory zero. The balance sheet can
be very large, long and short positions of equal size, but there is no exposure to the price \ts{10}{5}{4:57}. The rule comes from the view that banks, or their
dealer subsidiaries, took proprietary risk that, with FDIC insurance, fell on the taxpayer \ts{10}{5}{6:18}. It is about exposure to price risk, not leverage: ``a matched-book dealer can be
infinitely leveraged'' without price risk, and will try to persuade regulators it needs no capital. But there is no perfect hedge, and matched books carry other risks, liquidity risk among
them \ts{10}{5}{5:58}.

% ------------------------------------------------------------------
\section{Real-world dealers}\label{sec:dl-real}

Real dealers hedge: long some securities and short others that are correlated \ts{10}{6}{0:07}. Their balance sheet is the repo dealer's of lecture~7: reverse repo brings securities in (and
they may be sold short), repo sends them out as collateral for borrowing cash \ts{10}{6}{0:27}.

\begin{definition}[New concepts: gross and net exposure]
\emph{Gross exposure} is the size of the balance sheet; \emph{net exposure} is long positions minus short positions. A matched book can have a huge gross and zero net exposure; a small
one-sided book can have a small gross and a large net exposure \ts{10}{6}{1:50}.
\end{definition}

Typically (apart from the present distortions from QE3) dealers' net exposure is to the \emph{term spread}: long long-term bonds, short short-term ones, borrowing short and lending long
like a bank; they profit from the failure of the expectations hypothesis \ts{10}{6}{2:57}. The dealer business is building up gross exposure while controlling
net exposure \ts{10}{6}{3:59}. The Treynor model is about net exposure, price risk; gross exposure is liquidity risk, since the positions must be funded and counterparties may fail, a topic
for the banking lectures \ts{10}{6}{3:59}. Managing price risk well allows tighter spreads, which is how dealers compete \ts{10}{6}{4:41}.

\paragraph{How dealers make money} \ts{10}{6}{5:22}--\ts{10}{6}{7:48}:
\begin{enumerate}
  \item \textbf{Access to inventories}: knowing where securities are (which long-only pension funds hold what) and having a good relationship with a banker, perhaps the parent bank, for
        cash. The inventories are not on the dealer's balance sheet, where they would need funding, but ``out in the world'', like the mangoes in New Jersey;
  \item \textbf{watching the net position} and moving prices so as not to be ``bagged'' by better-informed traders, widening the spread when volatility rises, minute by minute;
  \item \textbf{offsetting} long with short exposures, buying at the bid and selling at the ask: the more both-way business, the more spread earned.
\end{enumerate}

\begin{intuition}[A just-in-time inventory system that straddles the hierarchy]
Through repo the dealer reaches the inventory of cash, through reverse the inventory of securities: it behaves as if it had inventories that are actually out in the market, a just-in-time
system. In a crisis it matters where the inventories are: cash comes from higher in the hierarchy, from banks, securities from lower, from security holders; the dealer straddles the
layers. Troubles come when it must produce securities it reversed in and sold, and more seriously when it must produce money it repoed in and spent \mn{10}{5}\mn{10}{6}.
\end{intuition}

\begin{coursenote}
The notes give the crisis examples: ``fails'' (collateral not returned) became so widespread that the Fed lent its own Treasuries to dealers through the Term Securities Lending Facility,
and when the repo market collapsed it opened the Primary Dealer Credit Facility, a lender of last resort for dealers \mn{10}{6}. (Both were created in March 2008, around the collapse of
Bear Stearns.)
\end{coursenote}

\begin{exercise}[Gross and net]
A dealer is long \$5 billion of 10-year Treasuries and short \$4.8 billion of 2-year Treasuries, all financed by repo and reverse repo.
\begin{enumerate}[(a)]
  \item Compute gross and net exposure in dollar terms.
  \item Why is the net exposure in dollars a poor measure of its price risk? What would you use instead?
\end{enumerate}
\end{exercise}

% ------------------------------------------------------------------
\section{Arbitrage and the assumption of perfect liquidity}\label{sec:dl-arbitrage}

A dealer long Treasury bonds and short Treasury bills is \emph{transporting liquidity} from one market to another: if every security needed its own dealer with its own capital, markets would
be much less liquid \ts{10}{7}{0:11}. Arbitrage links all markets: liquidity is a \emph{macroeconomic} concept, not market by market \ts{10}{7}{0:59}. Dealers choose
where to deploy their scarce resources, seeking wide spreads and little competition; so some markets have wide spreads, some steep bid--ask curves (trade in size and the price moves), some,
like Treasury bills, are nearly flat. Beyond this cross-section, macroeconomic fluctuations of liquidity move all prices together \ts{10}{7}{1:41}.

\paragraph{Prices away from value.} In Treynor's model the price moves with the dealer's \emph{inventory}, not with the fundamental value: the dealer lowers its price not because it knows the
asset is worth less but because its inventory is large \ts{10}{7}{3:05}. If fundamental value is the midpoint at zero inventory, then whenever the dealer holds
inventory, ``the price is wrong'' \ts{10}{7}{4:07}.

\begin{proposition}[Perfect liquidity is a limit case]
Most asset pricing, and general equilibrium theory, assume \emph{perfect arbitrage}: the outside spread collapses onto the fundamental value, value-based traders buy at any price a tiny bit below
it and sell a tiny bit above, the inside spread is squeezed to zero, and liquidity is a free good \ts{10}{7}{4:49}. In Treynor's world your
ability to trade comes from trading at a price \emph{different} from fundamental value: liquidity is provided for profit, and in fluctuating quantity \ts{10}{7}{6:36}.
\end{proposition}

The notes add a ``philosophical question'' \mn{10}{6}\mn{10}{7}: finance assumes that arbitrage creates perfect liquidity by entering at the smallest deviation; but in a fully efficient market
arbitrage would not be profitable, all bets being fair, so position takers would not make money, would compete less for market making, and markets would be less liquid, with wider spreads and
more volatile prices: ``in practice it seems that we must expect liquidity to enter into asset prices''. Arbitrageurs, by taking positions against temporary distortions, spread their
impact into other markets and so counteract them: ``porters'' of liquidity, producing a result ``no part of their intention''. Arbitrage and liquidity are two sides of the same coin
\mn{10}{6}. Some markets, and some times, are close to the frictionless limit; others are not \mn{10}{7}.

\begin{history}[Treynor, Black and ``Noise'']
Jack Treynor, a friend of Fischer Black, published ``The Economics of the Dealer Function'' in the \emph{Financial Analysts Journal} (1987); in 1971 he had already written, under the pseudonym
``Walter Bagehot'', ``The Only Game in Town'' on market makers and informed traders. Black, by then a partner at Goldman Sachs, gave his presidential address to the American Finance Association,
``Noise'', in New York in December 1985 (\emph{Journal of Finance}, July 1986). There he proposed to call a market efficient if price is within a factor of 2 of value, more than half and less than
twice value, and judged that by this definition almost all markets are efficient almost all of the time, ``almost all'' meaning at least 90\%. Mehrling reads the factor of 2 as Treynor's outside
spread \ts{10}{7}{6:58}. (In class the order of the two works is given the other way round.)\par
\emph{References:} J.~L.~Treynor, \emph{Financial Analysts Journal} 43(6), 1987; W.~Bagehot (pseud.), \emph{FAJ} 27(2), 1971; F.~Black, ``Noise'', \emph{Journal of Finance} 41(3), 1986.
\end{history}

\begin{reading}[Black, ``Noise'' (1986), section I ``Finance'', pp.~529--533]
\begin{itemize}
  \item Black contrasts trading on \emph{information} with trading on \emph{noise} as if it were information.
  \item Without noise traders there would be almost no trading in individual assets. An informed trader knows that only another informed trader takes the other side, and so declines to trade. Without trading in individual shares, portfolios, index futures and options could not be priced either.
  \item Noise trading makes markets liquid and makes it pay to collect costly information. But it also puts noise into prices: ``What's needed for a liquid market causes prices to be less efficient.''
  \item Informed traders do not eliminate the noise, because larger positions mean more risk and they can never be sure their information is not already in the price.
  \item Price wanders away from value like a drunk, and is pulled back the harder the further it strays. Hence the definition of efficiency: price within a factor of 2 of value.
\end{itemize}
The connection with Treynor: Black's noise traders are the flow the dealer absorbs, and his informed traders are Treynor's value-based traders. The band within which price may wander is the outside spread.
\end{reading}

``Money and banking matters'': it is not plumbing behind the scenes to be abstracted from in theories of value; it matters for asset pricing, income, employment and business cycles, as the crisis
showed \ts{10}{7}{9:03}.

% ------------------------------------------------------------------
\flushside
\section*{Key formulas and ideas}
\begin{keyformulas}
\begin{tabularx}{\linewidth}{@{}l L@{}}
Market vs funding liquidity & trade in volume without moving the price vs ability to raise funds\\
Inventories & the secret of continuous prices (supermarket, hierarchy of inventories)\\
Treynor & position limits (funding), outside spread (value-based traders), inside spread (width $\sim$ volatility, slope $\sim$ adverse selection)\\
Dealers & demand funding liquidity, supply market liquidity\\
Leverage & dealers trade on borrowed money; matched book: no price risk, but liquidity risk\\
Gross vs net & size of balance sheet vs longs $-$ shorts; dealers build gross, control net\\
Arbitrage & transports liquidity across markets: liquidity is macroeconomic\\
Perfect liquidity & outside spread $\to$ value, inside spread $\to$ 0: a logically impossible limit; Black: price within a factor 2 of value\\
\end{tabularx}
\end{keyformulas}

\section*{Exercises}

Standalone exercises follow the section they practise; connected or integrative ones are at the end.

\begin{exercise}[A linear dealer]
A dealer has $p^\ast=100$, inside spread $s=0.2$, position limits $\pm 1000$, and value-based bid and ask at 95 and 105.
\begin{enumerate}[(a)]
  \item Compute $\lambda$ and the quotes at inventories $-1000$, $0$, $500$, $1000$.
  \item The public sells the dealer 100 units, one at a time. What is the average price received, and how does it compare with $p^\ast$?
\end{enumerate}
\end{exercise}

\begin{exercise}[Crisis in the model]
In the model of the previous exercise the value-based bid falls to 70. Redraw the diagram. What happens to the dealer's willingness to absorb sales, and to the price path when the public sells
1500 units?
\end{exercise}
